\chapter{Synthese comparative}

\section{Cadre de comparaison}

Le coeur du projet consiste maintenant a comparer trois niveaux d'evaluation :
\begin{itemize}[leftmargin=*]
  \item les scores rapportes par les auteurs dans le papier ;
  \item notre baseline de reproduction, c'est-a-dire la version la plus fidele
  et stable du pipeline obtenue avec les ressources publiques ;
  \item notre meilleure configuration amelioree, issue des experiences
  d'amelioration conduites dans le chapitre precedent.
\end{itemize}

Cette structuration est importante, car elle permet de distinguer clairement :
\begin{itemize}[leftmargin=*]
  \item ce qui releve de la reproduction ;
  \item ce qui releve de la contribution personnelle ;
  \item ce qui reste encore ouvert a la fin du projet.
\end{itemize}

\section{Configurations retenues}

Deux configurations de reference emergent du travail realise.

\subsection{Baseline de reproduction}

La baseline de reproduction correspond a la configuration la plus proche du
papier que nous ayons pu executer de facon stable dans notre environnement :
\begin{itemize}[leftmargin=*]
  \item embeddings : \texttt{BAAI/bge-large-en-v1.5} ;
  \item evaluateur et namer : \texttt{mistralai/Mistral-7B-Instruct-v0.3} ;
  \item clusterer : \texttt{MiniBatchKMeans} ;
  \item \texttt{sample-size = 10} ;
  \item \texttt{sampler = random} ;
  \item \texttt{llm\_prompt\_style = simple} ;
  \item \texttt{naming\_sampler = head} ;
  \item \texttt{candidate\_k\_policy = fixed} ;
  \item merge final par labels uniquement.
\end{itemize}

\subsection{Meilleure configuration amelioree}

Au vu des ablations deja menees, la meilleure configuration amelioree retenue a
ce stade est la suivante :
\begin{itemize}[leftmargin=*]
  \item \texttt{llm\_prompt\_style = benchmark} ;
  \item \texttt{naming\_sampler = centroid} ;
  \item \texttt{candidate\_k\_policy = sqrt} ;
  \item merge final par labels uniquement.
\end{itemize}

Le merge \texttt{hybrid} et la politique \texttt{focused} pour \texttt{K} ne
sont donc pas retenus comme configurations finales, meme s'ils restent utiles
en tant qu'ablation et source d'enseignements sur le pipeline.

\section{Comparaison globale avec le papier}

Le tableau~\ref{tab:global-comparison} resume l'etat actuel du projet en
comparant, pour chaque benchmark, le score du papier, le score de reproduction
et le score de la configuration amelioree retenue a l'issue des ablations. Il
permet ainsi de lire directement ce que notre meilleure variante apporte par
rapport a la reproduction initiale, puis de situer ce resultat face au papier.

\begin{table}[h]
  \centering
  \small
  \begin{tabular}{lrrrrrl}
    \toprule
    Benchmark & Papier & Reproduction & Amelioree & Gain vs reprod. & Delta vs papier & Meilleure config \\
    \midrule
    Bank77 & 82.32 & 83.78 & 82.28 & -1.50 & -0.04 & reproduction \\
    CLINC(I) & 94.12 & 90.90 & 92.36 & +1.46 & -1.76 & amelioree \\
    MTOP(I) & 72.45 & 71.12 & 73.08 & +1.96 & +0.63 & amelioree \\
    Massive(I) & 78.12 & 72.44 & 75.68 & +3.24 & -2.44 & amelioree \\
    \bottomrule
  \end{tabular}
  \caption{Comparaison globale entre le papier, la reproduction et la configuration amelioree finale.}
  \label{tab:global-comparison}
\end{table}

Ce tableau montre deja une image assez nette de l'etat du projet :
\begin{itemize}[leftmargin=*]
  \item la configuration amelioree n'est pas universellement meilleure :
  Bank77 se degrade par rapport a la reproduction, meme s'il reste quasiment
  aligne sur le score du papier ;
  \item les ameliorations sont en revanche benefiques sur CLINC(I), MTOP(I) et
  Massive(I), avec des gains respectifs de 1.46, 1.96 et 3.24 point de NMI par
  rapport a la reproduction ;
  \item MTOP(I) devient particulierement interessant, car la version amelioree
  depasse legerement la cible du papier.
\end{itemize}

\section{Comparaison detaillee sur les benchmarks ameliores}

Les ameliorations n'ont pas toutes ete lancees sous forme d'ablations
exploratoires sur les quatre benchmarks. En l'etat, la comparaison la plus
complete entre variantes concerne CLINC(I) et Massive(I), car ce sont les deux
benchmarks sur lesquels les ecarts de reproduction etaient les plus
significatifs.

Le tableau~\ref{tab:improvement-path} resume la trajectoire observee.

\begin{table}[h]
  \centering
  \small
  \begin{tabular}{lrrrrrr}
    \toprule
    Benchmark & Reprod. & Imp. 1 & Imp. 2 & Imp. 3 & Imp. 4 & Papier \\
    \midrule
    CLINC(I) & 90.90 & 91.79 & 91.44 & 92.36 & 92.01 & 94.12 \\
    Massive(I) & 72.44 & 73.77 & 73.69 & 75.68 & 73.82 & 78.12 \\
    \bottomrule
  \end{tabular}
  \caption{Trajectoire des scores au fil des ameliorations. La colonne Reprod. reprend la baseline officielle de la campagne finale ; les colonnes d'amelioration correspondent aux meilleures variantes exploratoires observees pour chaque etape. Imp. 1 : coherence / naming. Imp. 2 : merge hybride. Imp. 3 : \texttt{K} adaptatif \texttt{sqrt}. Imp. 4 : \texttt{K} \texttt{focused}.}
  \label{tab:improvement-path}
\end{table}

Cette trajectoire permet de tirer plusieurs conclusions robustes :
\begin{itemize}[leftmargin=*]
  \item l'amelioration du bloc \texttt{coherence / naming} apporte un premier
  gain utile sur CLINC(I), mais reste instable sur Massive(I) ;
  \item le merge hybride aide a la consolidation, sans devenir la meilleure
  configuration globale ;
  \item la politique \texttt{sqrt} pour \texttt{K} est l'amelioration la plus
  forte et la plus stable parmi celles testees ;
  \item la politique \texttt{focused} ne depasse pas \texttt{sqrt}, mais a mis
  en evidence un besoin de robustesse supplementaire dans l'evaluateur local.
\end{itemize}

\section{Lecture scientifique de l'etat actuel}

Le resultat principal du projet peut se formuler ainsi : nous avons obtenu une
reproduction credible et defendable du papier, puis nous avons montre que des
ameliorations ciblees du pipeline peuvent reduire une partie de l'ecart au
papier sur plusieurs benchmarks, sans pour autant produire un gain uniforme.

Le message scientifique n'est donc ni :
\begin{itemize}[leftmargin=*]
  \item ``nous avons reproduit exactement le papier'' ;
  \item ni ``nous avons simplement teste quelques variantes sans ligne
  directrice''.
\end{itemize}

Le message plus juste est le suivant :
\begin{itemize}[leftmargin=*]
  \item nous avons reconstruit le coeur de la methode dans une forme runnable ;
  \item nous avons mesure objectivement les ecarts a la version des auteurs ;
  \item nous avons identifie les leviers experimentaux les plus importants ;
  \item nous avons obtenu une meilleure configuration que notre reproduction de
  depart sur trois benchmarks sur quatre, avec un comportement plus nuance sur
  Bank77.
\end{itemize}

\section{Etat courant du projet}

A la date de ce rapport, la meilleure base experimentale du projet est donc :
\begin{itemize}[leftmargin=*]
  \item une baseline de reproduction stable sur les quatre benchmarks ;
  \item une baseline amelioree uniforme, basee sur \texttt{prompt = benchmark},
  \texttt{naming = centroid} et \texttt{candidate\_k\_policy = sqrt} ;
  \item un tableau comparatif complet qui permet de trancher, benchmark par
  benchmark, entre reproduction et version amelioree.
\end{itemize}

La derniere etape logique n'est donc plus experimentale, mais redactionnelle :
il s'agit maintenant d'exploiter proprement cette comparaison finale dans le
rapport, puis de la replacer dans un cadre plus large avec l'etat de l'art et
la discussion generale du projet.
